% TOP_PROBLEM: {"id": 2990, "title": "Kirby Problem 4.114", "category_id": 7, "category": {"id": 7, "name": "topology", "display_name": "Topology", "description": "Properties preserved under continuous deformations.", "slug": "topology", "order_index": 7, "created_at": "2026-07-31T15:26:25.671Z"}, "status": "open"}
% FIRST_POSTED: null
% ENTRY_KIND: statement_only
% Imported from: batch14.tex; UnsolvedMath ID 2990
\documentclass[11pt]{article}
\usepackage[margin=25mm]{geometry}
\usepackage{fontspec}
\setmainfont{Latin Modern Roman}
\usepackage{amsmath,amssymb,mathtools}
\usepackage{unicode-math}
\setmathfont{Latin Modern Math}
\usepackage{xurl,hyperref}
\hypersetup{colorlinks=true,urlcolor=blue}
\setcounter{secnumdepth}{0}
\setlength{\parindent}{0pt}
\setlength{\parskip}{5pt}
\setlength{\emergencystretch}{3em}
\newcommand{\sref}[2]{\hyperlink{source:#1:#2}{[S#2]}}
\newcommand{\cataloguescope}[1]{\paragraph{Recorded target.}#1}
\begin{document}
% UnsolvedMath ID 2990; source catalogue code KP-4.114
\setcounter{section}{2989}
\section{Kirby Problem 4.114}\label{2990}
{\small\sffamily UnsolvedMath ID 2990 · Topology}
\subsection{Definitions and mathematical statement}

\paragraph{Shared notation.}$\mathbb N=\{1,2,\ldots\}$, and $\mathbb Z,\mathbb Q,\mathbb R,\mathbb C$ denote the integers, rationals, reals and complex numbers. Any other conventions are specified locally.


A Heegaard splitting of a closed orientable smooth three-manifold $Y$ is $Y=H_-\cup_\Sigma H_+$ with $H_\pm$ handlebodies. A Morse function $f:Y\to\mathbb R$ is self-indexing if the Hessian at every critical point is nonsingular and the critical value is its Morse index (the number of negative Hessian eigenvalues). The level $f^{-1}(3/2)$ separates the index-zero/one handles from the index-two/three handles and induces the splitting.

For smooth maps $f_i:M_i\to N_i$, right-left equivalence means that there exist diffeomorphisms $\phi:M_0\to M_1$ and $\psi:N_0\to N_1$ with $f_1=\psi\circ f_0\circ\phi^{-1}$. If $N_0=N_1$ and $\psi$ is the identity, this is right equivalence.
A genus-$g$ trisection of a closed connected oriented smooth four-manifold $X$ is a decomposition $X=X_1\cup X_2\cup X_3$ such that $X_i$ is a four-dimensional $1$-handlebody (a boundary connected sum of $k_i$ copies of $S^1\times B^3$), each $X_i\cap X_j$ is a three-dimensional genus-$g$ handlebody, and $X_1\cap X_2\cap X_3=\Sigma_g$ is a closed genus-$g$ surface. Corners in these descriptions are rounded.

A trisected Morse $2$-function is a generic smooth map to a disk with only fold and cusp singularities, whose inverse images of three sectors give a trisection. In local smooth coordinates the fold and cusp forms are respectively
\[
 (u,v,w,z)\mapsto(u,\ \pm v^2\pm w^2\pm z^2),\qquad
 (u,v,w,z)\mapsto(u,\ v^3+uv\pm w^2\pm z^2).
\]
The standard trisected arrangement has one definite fold circle at the boundary, radial rays crossing the indefinite folds in the index-two direction, and three sectors containing at most one cusp on each fold arc. A general trisection may have unequal parameters $k_1,k_2,k_3$; equality is not imposed. A simplified trisection is induced by such a map with embedded singular-value curves and with cusps occurring in triples on innermost fold circles; crossings of indefinite folds between the sectors are excluded. \sref{2990}{3}
\paragraph{Statement recorded in the supplied source.}
For a prescribed Heegaard splitting of $Y$, classify the self-indexing Morse functions inducing that splitting under right-left equivalence, and under right equivalence. For a prescribed trisection of a closed orientable smooth four-manifold $X$, classify the generic trisected Morse $2$-functions $X\to\mathbb R^2$ inducing that trisection under each of the same equivalences. Include the subclass inducing a prescribed simplified trisection. In particular, determine whether equivalent induced trisections force right-left equivalent maps, and describe the failure if they do not. \sref{2990}{2}
\subsection{Short English statement}
Classify all Morse functions or generic planar maps that realize one fixed Heegaard splitting or trisection, with or without changing the target coordinates.
\subsection{Sources}
\begin{description}
\item[\hypertarget{source:2990:1}{S1}] ULAM.AI and the UnsolvedMath Contributors, \emph{UnsolvedMath: A Curated Collection of Open Mathematics Problems}, record 2990 (KP-4.114). CC BY 4.0. \url{https://huggingface.co/datasets/ulamai/UnsolvedMath}.
\item[\hypertarget{source:2990:2}{S2}] R. İnanç Baykur, Robion C. Kirby and Daniel Ruberman (eds.), \emph{K3: A New Problem List in Low-Dimensional Topology}, Mathematical Surveys and Monographs 295, American Mathematical Society (2026), Problem 4.114, author preliminary version. \url{https://math.berkeley.edu/sites/default/files/surv-295-ruberman-watermarked-author-pdf.pdf}.
\item[\hypertarget{source:2990:3}{S3}] R. İnanç Baykur and Osamu Saeki, \emph{Simplified broken Lefschetz fibrations and trisections of 4-manifolds}, Proceedings of the National Academy of Sciences 115 (2018), 10894--10900, §§2--3 (arXiv:1710.06529v1, 2017). \url{https://arxiv.org/abs/1710.06529}.
\end{description}
\label{end:2990}
\end{document}
